\section{A reproducible serving experiment}\label{sec:C31:08}
The experimental unit is a trace-defined session with declared model, tokenizer, harness, and memory behavior. Record input and output lengths, reused prefix, cache tier, context edits, tool waits, retries, and termination. Run both open-loop offered-load tests and closed-loop concurrency tests, because closed-loop clients reduce arrivals when latency increases and can hide overload. Separate cold starts, warm kernels, cold caches, warm prefixes, and resumed sessions on different replicas.

Report time to first token, inter-token latency, request completion, session completion, throughput, rejection, timeout, and failure. Provide p50, p95, and p99 where sample size supports them, with the number of observations and a clear treatment of censored requests. A percentile from completed requests alone can look better when the slowest requests fail. Include tool time and memory maintenance in task-level results even when engine-only timing is also reported.

The ablation sequence should isolate one mechanism at a time: no prefix reuse versus reuse, uniform versus hybrid allocation, cache offload disabled versus enabled, fixed versus adaptive speculative budget, and monolithic versus disaggregated execution. Hold the action-output quality criterion fixed. If an optimization changes context policy or truncates evidence, report that as a behavioral intervention rather than an invisible serving improvement.

A negative result is informative. Tiered caching may lose on low-reuse traffic; speculation may lose with weak acceptance; a larger batch may improve aggregate throughput while delaying critical tool actions; compressed KV may preserve average language quality while corrupting rare arguments. The deployment decision follows from the joint quality, latency, cost, and reliability envelope. Peak throughput is one coordinate of that envelope, not a sufficient description of an agent-serving system.
